\documentclass[master=ucll,11pt,dutch,twoside]{kulemt}
%http://texdoc.net/texmf-dist/doc/latex/memoir/memman.pdf

%\documentclass[master=ucll,11pt,dutch,twoside]{ucllfilip}

%master is hetgene dat gedefinieerd staat in de .cfg file
%masteroption is hetgene bij ELT gedeelte staat als optie
%and command is voor een 2de lijn bij te voegen
%%  author={Filip Vanden Eynde \and R0363898 \and 3TX/1},
%  assistant={Filip Assistent},
 %assessor={Filip Assessor},
\setup{
	title={Samenvatting IT en RECHT 2017-2018},
  author={Filip Vanden Eynde \and R0363898 \and 3TX/1},
  teacher={Jan-Pieter Everts},
  company={UCLL Leuven-Limburg},
  location={Haasrode, Leuven, België},
  textsubject={Samenvatting voor het vak IT en RECHT in het academiejaar 2017-2018 voor de Professionele Bachelor
  Toegepaste Informatica},
  keywordspdf={Samenvatting, IT en Recht, Filip Vanden Eynde, schooljaar 2017-2018, UCLL}
  }
  
\newcommand{\quotes}[1]{``#1''}
  
% De volgende \setup mag verwijderd worden als geen fiche gewenst is.
%\setup{filingcard,
 % translatedtitle={The best master's thesis ever},
 % udc=621.3,
  %keywords={Stageverslag, Dunedin, Nieuw-Zeeland, UCLL, schooljaar 2016-2017, Filip Vanden Eynde},
  %shortabstract={Hier komt een heel bondig abstract van hooguit 500
   % woorden. \LaTeX\ commando's mogen hier gebruikt worden. Blanco lijnen
   % (of het commando \texttt{\string\pa r}) zijn wel niet toegelaten!
   % \endgraf \lipsum[2]}}
	
% Verwijder de "%" op de volgende lijn als je de kaft wil afdrukken
%\setup{coverpageonly}
%\write18{pdflatex masterproefcoverpage -aux-directory=outputfoldercoverpage}
%\write18{pdflatex masterproefcoverpage}

% Kies de fonts voor de gewone tekst, bv. Latin Modern
%\setup{font=palatino}

\setup{font=utopia}

\setup{inputenc=utf8}

\usepackage{ifthen}
\newboolean{pdfoutput}
\setboolean{pdfoutput}{true}
%\setboolean{pdfoutput}{false}

% Hier kun je dan nog andere pakketten laden of eigen definities voorzien

\usepackage{kulemtx}
\chapterstyle{kulemtman}
\kulemtmanToC{}

\usepackage{booktabs}
\usepackage{colortbl}
\usepackage[table,dvipsnames]{xcolor}
\usepackage{subcaption}


%\begin{comment}

%for bibliography
\usepackage[backend=biber,
%%style=alphabetic,
style=numeric,
%%style=ieee,
%%style=ieee-alphabetic,
citestyle=numeric,
%%refsection=chapter,
%refsegment=chapter,
%defernumbers=true,
%%sorting=ynt,
sorting=nty,
sortcites, % Dit zorgt ervoor dat als je meerdere referenties toont zoals [9,1] dat dit wordt getoond als [1,9] dus numeriek gesorteerd
]{biblatex}
\addbibresource{referenties.bib}

%\end{comment}


%indexes
%voor indexen te maken
\usepackage{imakeidx}
\makeindex[
%columns=3, % geeft layout van 3 kolommen
title=Alphabetical Index, % verandert de default title = index
intoc] % voegt de indexpagina aan de inhoudsopgave toe

\usepackage{import}
\usepackage{enumitem}
\usepackage{pdfpages}
\usepackage{todonotes}

%provides the multicols environment which typesets text into multiple columns.
\usepackage{multicol}
\usepackage[official]{eurosym}

%-------------------------------------------------------------------------------------------------------------------------
% pas plaats voor en na headings aan gebruikmakend van package titlesec
%-------------------------------------------------------------------------------------------------------------------------

\usepackage{titlesec}
%\titlespacing\section{0pt}{12pt plus 4pt minus 2pt}{0pt plus 2pt minus 2pt}
%\titlespacing\subsection{0pt}{12pt plus 4pt minus 2pt}{0pt plus 2pt minus 2pt}

%------------------------------------------------------------------------

%hier komen alle commando's en constante variabelen voor in de opdracht te gebruiken.

\graphicspath{{images/}} % Set the default folder for images
%I usually write something like this, so the entry is hyperlinked for onscreen reading but there's also a footnote to the URL for paper output.
\newcommand\fnurl[2]{%
	\href{#2}{#1}\footnote{\url{#2}}%
}
%\setlist[itemize,1]{leftmargin=*,topsep=0pt}
%\setlist[enumerate,1]{leftmargin=*,topsep=0pt}
\setlist{nosep}

\renewcommand\mempostaddapppagetotochook{\cftinserthook{toc}{BREAK}}
\cftinsertcode{BREAK}{\changetocdepth{-10}}
\let\normalchangetocdepth\changetocdepth % needed for later

\makeatletter
\newcommand\appendixtableofcontents{
	\begingroup
	\let\changetocdepth\@gobble
	\normalchangetocdepth{-10}
	\cftinsertcode{BREAK}{\normalchangetocdepth{3}}
	\renewcommand\contentsname{Appendices overzicht}
	\tableofcontents*
	\endgroup
}
\makeatother


\ifthenelse {\boolean{pdfoutput}}
	{
	\usepackage[
	%hidelinks = true,
	%hidelinks=\tableofcontents,

    pdfusetitle, colorlinks, urlcolor=blue, linkcolor=blue,
	plainpages=false ,bookmarks=true, bookmarksnumbered, linktoc=all, pdffitwindow=true,
	pdftex,pdfauthor={\theauthor},pdftitle={\thetitle},pdfsubject={\thesubject},
	pdfkeywords={\thekeywordspdf},pdfproducer={\theauthor},pdfcreator={\theauthor}
	]{hyperref}	
	}{

\usepackage[
    pdfusetitle,
    colorlinks,
    urlcolor=black,
    linkcolor=black,
    citecolor=black,
	plainpages=false,
	bookmarks=true,
	bookmarksnumbered,
	linktoc=all,
	pdffitwindow=true,
	pdftex,
	pdfauthor={\theauthor},
	pdftitle={\thetitle},
	pdfsubject={\thesubject},
	pdfkeywords={\thekeywordspdf},
	pdfproducer={\theauthor},
	pdfcreator={\theauthor}
	]{hyperref}	
}

%for plural forms of acronyms
\usepackage{acronym}

\usepackage{xparse}
\DeclareDocumentCommand{\newdualentry}{ O{} O{} m m m m } {
	\newglossaryentry{gls-#3}{name={#5},text={#5\glsadd{#3}},
		description={#6},#1
	}
	\makeglossaries
	\newacronym[see={[Glossary:]{gls-#3}},#2]{#3}{#4}{#5\glsadd{gls-#3}}
}

\usepackage[acronym, toc]{glossaries}
\usepackage{glossary-mcols}
\makeglossaries

%A
\begin{comment}
%==> dit is het woord waarnaar verwezen wordt in de tekst
\newglossaryentry{}{
%==> dit is hetgene dat wordt getoond in de glossaries lijst
name={},
%==> dit is de beschrijving van de glossary
description={}
}

%==> dit is het woord waarnaar verwezen wordt in de tekst
\newglossaryentry{}{
%==> dit is hetgene dat wordt getoond in de glossaries lijst
name={},
%==> dit is de beschrijving van de glossary
description={}
}

%==> dit is het woord waarnaar verwezen wordt in de tekst
\newglossaryentry{}{
%==> dit is hetgene dat wordt getoond in de glossaries lijst
name={},
%==> dit is de beschrijving van de glossary
description={}
}

%==> dit is het woord waarnaar verwezen wordt in de tekst
\newglossaryentry{}{
%==> dit is hetgene dat wordt getoond in de glossaries lijst
name={},
%==> dit is de beschrijving van de glossary
description={}
}

%==> dit is het woord waarnaar verwezen wordt in de tekst
\newglossaryentry{}{
%==> dit is hetgene dat wordt getoond in de glossaries lijst
name={},
%==> dit is de beschrijving van de glossary
description={}
}


\end{comment}

\newglossaryentry{accesscontrollist}{
name={Access Control List },
description={ }
}

\newglossaryentry{latex}
{
	name=latex,
	description={Is a mark up language specially suited 
		for scientific documents}
}

\newglossaryentry{maths}
{
	name=mathematics,
	description={Mathematics is what mathematicians do}
}
% label for reference -> acronym itself -> explanation
%\newacronym{ } { } {}

% label for reference -> acronym itself -> explanation
\newacronym{acl}{ACL}{Access Control List}

\newacronym{gcd}{GCD}{Greatest Common Divisor}

\newacronym{lcm}{LCM}{Least Common Multiple}
%\import{./}{glossary.tex}
%\import{./}{acronyms.tex}
%\import{glossary}

\usepackage{sty/renewcommandsfilip}


%%%%%%%
% Om wat tekst te genereren wordt hier het lipsum pakket gebruikt.
% Bij een echte masterproef heb je dit natuurlijk nooit nodig!
\IfFileExists{lipsum.sty}%
 {\usepackage{lipsum}\setlipsumdefault{11-13}}%
 {\newcommand{\lipsum}[1][11-13]{\par Hier komt wat tekst: lipsum ##1.\par}}
%%%%%%%

%\includeonly{hfdst-n}
\begin{document}

\begin{preface}
  Dit is mijn dankwoord om iedereen te danken die mij bezig gehouden heeft.
  Hierbij dank ik mijn promotor, mijn begeleider en de voltallige jury.
  Ook mijn familie heeft mij erg gesteund natuurlijk.
  %\quotes{Otago Polytechnic}
  %stagebegeleider \stagebegeleider
\end{preface}

\tableofcontents*


\begin{abstract}
  In dit \texttt{abstract} environment wordt een al dan niet uitgebreide
  samenvatting van het werk gegeven. De bedoeling is wel dat dit tot
  1~bladzijde beperkt blijft.

  \lipsum[1]
    \todo[inline]{Samenvatting schrijven}

\end{abstract}

% Een lijst van figuren en tabellen is optioneel
%\listoffigures
%\listoftables
% Bij een beperkt aantal figuren en tabellen gebruik je liever het volgende:
\listoffiguresandtables
%\todo[inline]{Nakijken of alle labels van figuren en tabellen uniek en duidelijk gedefinieerd zijn?}

% De lijst van symbolen is eveneens optioneel.
% Deze lijst moet wel manueel aangemaakt worden, bv. als volgt:
\chapter{Lijst van afkortingen en symbolen}
\section*{Afkortingen}
\begin{flushleft}
  \renewcommand{\arraystretch}{1.1}
  \begin{tabularx}{\textwidth}{@{}p{12mm}X@{}}
    LoG   & Laplacian-of-Gaussian \\
    MSE   & Mean Square error \\
    PSNR  & Peak Signal-to-Noise ratio \\
  \end{tabularx}
\end{flushleft}
\section*{Symbolen}
\begin{flushleft}
  \renewcommand{\arraystretch}{1.1}
  \begin{tabularx}{\textwidth}{@{}p{12mm}X@{}}
    42    & ``The Answer to the Ultimate Question of Life, the Universe,
            and Everything'' volgens de \cite{h2g2} \\
    $c$   & Lichtsnelheid \\
    $E$   & Energie \\
    $m$   & Massa \\
    $\pi$ & Het getal pi \\
  \end{tabularx}
\end{flushleft}

%\todo[inline]{Op het einde de list of todos verwijderen.}
%\listoftodos

% Nu begint de eigenlijke tekst
\mainmatter

%allemaal new stuff
%\clearforchapter

%end allemaal new stuff

%----------------------------------------

%Dit betekent pak de pagestyle voor de eerste pagina als een nieuw hoofdstuk begint
\aliaspagestyle{chapter}{chapternohead}
%\pagestyle{chapternohead}
\pagestyle{ruledfilip}

%origineel
%\chapter{Inleiding}
\label{inleiding}
%\thispagestyle{chapternohead}
%\pagestyle{ruledfilip}
In dit hoofdstuk wordt het werk ingeleid. Het doel wordt gedefinieerd en er
wordt uitgelegd wat de te volgen weg is (beter bekend als de rode draad).

Als je niet goed weet wat een masterproef is, kan je altijd
Wikipedia\cite{wiki} eens nakijken.

\section{Lorem ipsum 4--5}
\lipsum[4-5]

\section{Lorem ipsum 6--7}
\lipsum[6-7]



%%% Local Variables: 
%%% mode: latex
%%% TeX-master: "masterproef"
%%% End: 


\part{Sociaal Recht}

%\import{texfiles/include-chapter1/}{chapter1.tex}

\chapter{Algemene bepalingen van de arbeidsovereenkomstenwet van          3 juli 1978}
\label{hoofdstuk:1}

%\setcounter{section}{-1}

\section{Toepassingsgebied van de Arbeidsovereenkomst}

\subsection{Wat is een arbeidsovereenkomst?}

\subsection{Het ondergeschikte verband}

\subsubsection{Algemeen}

\subsubsection{Schijnzelfstandigen - Pseudowerknemers}

\subsubsection{Samengaan van een arbeidsovereenkomst met een aannemings- of zelfstandigenovereenkomst}

\subsection{Toepassingsgebied van de Arbeidsovereenkomstenwet}

\section{Het aangaan van de arbeidsovereenkomst}

\subsection{Geldigheidsvereisten}

\subsubsection{Bekwaamheid}

\subsubsection{Toestemming}

\paragraph{Geweld}

\paragraph{Dwaling}

\paragraph{Bedrog}

\subsubsection{Voorwerp en oorzaak}

\subsubsection{Gevolgen}

\subsection{Vormvereisten van de arbeidsovereenkomst}

\section{Soorten arbeidsovereenkomsten}

\subsection{Soorten naargelang de beoogde arbeid}

\subsubsection{Werklieden}

\subsubsection{Bedienden}

\subsubsection{Handelsvertegenwoordigers}

\subsubsection{Dienstboden}

\subsubsection{Studenten}

\subsubsection{Huisarbeid}

\subsection{Soorten naargelang de duur van de overeenkomst}

\subsubsection{De arbeidsovereenkomst voor onbepaalde tijd}

\subsubsection{De arbeidsovereenkomst voor een bepaalde tijd}

\paragraph{Onderscheid tussen overeenkomsten van onbepaalde en bepaalde tijd}

\paragraph{Opeenvolgende arbeidsovereenkomsten voor een bepaalde tijd}

\paragraph{Afwijking van het principiële verbod om opeenvolgende arbeidsovereenkomsten voor een bepaalde tijd te sluiten}

\paragraph{Algemeen beginsel van non-discriminatie ten gunste van werknemers met een overeenkomst voor bepaalde tijd}

\subsubsection{De arbeidsovereenkomst voor duidelijk omschreven werk}

\subsubsection{De vervangingsovereenkomst}

\subsection{Soorten naargelang de omvang van het werk}

\subsubsection{Voltijdse arbeidsovereenkomst}

\subsubsection{Arbeidsovereenkomst voor deeltijdse arbeid}

\paragraph{Geschrift}

\paragraph{De arbeidsregeling}

\paragraph{Het werkrooster}

\paragraph{Minimumgrenen van de arbeidsduur}

\paragraph{Bekendmaking van de variabele werkroosters}

\paragraph{Voorrang bij het overstappen naar een voltijdse betrekking}

\paragraph{Non-discriminatie ten gunste van deeltijdse werknemers}

\paragraph{Rechten deeltijdse werknemer versus voltijdse werknemer}


\subsection{Soorten naargelang de wijze waarop de overeenkomst werd gesloten}

\subsubsection{Mondelinge en schriftelijke arbeidsovereenkomsten}

\subsubsection{Elektronische arbeidsovereenkomsten}

%%% Local Variables: 
%%% mode: latex
%%% TeX-master: "masterproef"
%%% End: 

\chapter{Inhoud van de arbeidsovereenkomst}
\label{hoofdstuk:2}

\section{Het dwingende karakter van de Arbeidsovereenkomst}

\subsection{Dwingend recht, suppletoire betalingen en openbare orde}

\subsection{Artikel 6 van de Arbeidsovereenkomst}

\section{Het proefbeding}

\section{Het concurrentiebeding}

\subsection{Concurrentie}

\subsection{Wat is een concurrentiebeding?}

\subsection{Gewoon concurrentiebeding voor arbeiders en bedienden}

\subsubsection{Jaarloongrens}

\subsubsection{Beperking tot soortgelijke activiteiten}

\subsubsection{Geografische beperking}

\subsubsection{Beperking in de tijd}

\subsubsection{Forfaitaire vergoeding}

\subsubsection{Vormvereisten}

\subsubsection{Sanctie}

\subsubsection{Wanneer moet een concurrentiebeding toegepast worden?}

\subsubsection{Schadevergoeding bij overtreding van het concurrentiebeding}

\subsection{Afwijkend concurrentiebeding voor sommige bedienden}

\subsection{Het concurrentiebeding voor handelsvertegenwoordigers}

\section{Het scholingsbeding}

\subsection{Wat is een scholingsbeding?}

\subsection{Voorwaarden voor een geldig scholingsbeding}

%%% Local Variables: 
%%% mode: latex
%%% TeX-master: "masterproef"
%%% End: 

\chapter{Uitvoering van de arbeidsovereenkomst}
\label{hoofdstuk:2}

\section{Verplichtingen van beide partijen}

\subsection{De goede trouw in de uitvoering van de arbeidsovereenkomst}

\subsection{Wederzijdse eerbied en achting}

\subsection{Psychosociale risico's op het werk}

\subsection{Privacy van de werknemer versus controlerecht van de werkgever}

\subsubsection{Beperking van het controlerecht door de werkgever}

\subsubsection{Waarde van onrechtmatig verkregen bewijsmateriaal}

\section{Verplichtingen van de werknemer}

\subsection{Arbeid verrichten}

\subsection{Bevelen en instructies}

\subsection{Veiligheid}

\subsection{Arbeidsgereedschap}

\subsection{Concurrentie}

\section{Verplichtingen van de werkgever}

\subsection{Arbeid verschaffen}

\subsection{Hulp}

\subsection{Veiligheid en gezondheid}

\subsection{Loon}

\subsection{Huisvesting en voeding}

\subsection{Geloofsplichten en burgerlijke plichten}

\subsection{Onthaal}

\subsection{Arbeidsgereedschap}

\subsection{Documenten}

\section{Aansprakelijkheid van de partijen}

\subsection{Aansprakelijkheid van de werkgever}

\subsection{Aansprakelijkheid van de werknemer}



%%% Local Variables: 
%%% mode: latex
%%% TeX-master: "masterproef"
%%% End: 

\chapter{Schorsing van de uitvoering van de arbeidsovereenkomst}
\label{hoofdstuk:2}

\section{Het begrip schorsing}

\section{Overmacht}

\subsection{Begrip}

\subsection{Gevolgen}

\subsection{Recht op werkloosheidsuitkeringen}

\section{Het gewaarborgd dagloon}

\subsection{Met vertraging of niet op het werk aankomen}

\subsection{De arbeid niet kunnen beginnen of voortzetten}

\section{Arbeidsongeschiktheid}

\subsection{Schorsing}

\subsection{Verplichtingen van de werknemer}

\subsubsection{Meldingsplicht}

\subsubsection{Geneeskundig getuigschrift}

\subsection{Controle door de werkgever}

\subsection{Betwistingen van medische aard}

\subsection{Gewaarborgd loon}

\subsubsection{Algemeen}

\subsubsection{Gewaarborgd dagloon (art. 27 WAO) en gewaarborgd loon bij arbeidsongeschiktheid}

\subsubsection{Uitsluiting}

\subsubsection{Recht van verhaal tegen de aansprakelijke derde}

\subsection{Het bedrag van het gewaarborgd loon}

\subsubsection{De regeling voor werklieden (WAO en cao nr. 12\textit{bis}}

\paragraph{Carenzdag}

\paragraph{Periode en bedrag}

\paragraph{Anciënniteitsvoorwaarde}

\paragraph{Hervallen}

\paragraph{Gewaarborgd loon bij arbeidsongeval of beroepsziekte}

\subsubsection{De regeling voor bedienden}

\paragraph{Algemeen stelsel}

\paragraph{Bedienden die de regeling voor werklieden volgen}

\paragraph{Gewaarborgd loon bij arbeidsongeval of beroepsziekte}

\subsubsection{Overzicht van het gewaarborgd loon bij ziekte}

\section{Kort verzuim (of ook: klein verlet)}

\subsection{Algemene principes}

\subsection{Lijst van de gebeurtenissen}

\subsection{Geboorteverlof}

\subsection{Adoptieverlof}



%%% Local Variables: 
%%% mode: latex
%%% TeX-master: "masterproef"
%%% End: 

\chapter{Een volgend hoofdstuk}
\label{hoofdstuk:2}
Een hoofdstuk behandelt een samenhangend geheel dat min of meer op zichzelf
staat. Het is dan ook logisch dat het begint met een inleiding, namelijk
het gedeelte van de tekst dat je nu aan het lezen bent.

\section{Eerste onderwerp in dit hoofdstuk}
De inleidende informatie van dit onderwerp.

\subsection{Een item}
Een tekst staat nooit alleen. Dit wil zeggen dat er zeker ook referenties
nodig zijn. Dit kan zowel naar on-line documenten\cite{wiki} als naar
boeken\cite{pratchett06:_good_omens}.

\section{Figuren}
Figuren worden gebruikt om illustraties toe te voegen. Dit is dan ook de
manier om beeldmateriaal toe te voegen zoals getoond wordt in
figuur~\ref{fig:logo}.

\begin{figure}
  \centering
  \includegraphics{logokul}
  \caption{Het KU~Leuven logo.}
  \label{fig:logo}
\end{figure}

\section{Tabellen}
Tabellen kunnen gebruikt worden om informatie op een overzichtelijke te
groeperen. Een tabel is echter geen rekenblad! Vergelijk maar eens
tabel~\ref{tab:verkeerd} en tabel~\ref{tab:juist}. Welke tabel vind jij het
duidelijkst?

\begin{table}
  \centering
  \begin{tabular}{||l|lr||} \hline
    gnats     & gram      & \$13.65 \\ \cline{2-3}
              & each      & .01 \\ \hline
    gnu       & stuffed   & 92.50 \\ \cline{1-1} \cline{3-3}
    emu       &           & 33.33 \\ \hline
    armadillo & frozen    & 8.99 \\ \hline
  \end{tabular}
  \caption{Een tabel zoals het niet moet.}
  \label{tab:verkeerd}
\end{table}

\begin{table}
  \centering
  \begin{tabular}{@{}llr@{}} \toprule
    \multicolumn{2}{c}{Item} \\ \cmidrule(r){1-2}
    Animal    & Description & Price (\$)\\ \midrule
    Gnat      & per gram    & 13.65 \\
              & each        & 0.01 \\
    Gnu       & stuffed     & 92.50 \\
    Emu       & stuffed     & 33.33 \\
    Armadillo & frozen      & 8.99 \\ \bottomrule
  \end{tabular}
  \caption{Een tabel zoals het beter is.}
  \label{tab:juist}
\end{table}

\section{Lorem ipsum}
Tenslotte gaan we hier nog wat tekst voorzien zodat er minstens een
bijkomende bladzijde aangemaakt wordt. Dat geeft de gelegenheid om eens te
zien hoe de koptekst en de voettekst zich gedragen.

\subsection{Lorem ipsum dolor sit amet, consectetur adipiscing elit}
Sed nec tortor id felis tristique sodales. Nulla nec massa eu dui fermentum
tincidunt. Integer ullamcorper ante eget eros posuere faucibus. Nam id
ligula ut augue pulvinar vulputate id at purus. Aenean condimentum tortor
eu mi placerat eget eleifend massa mollis. Nam est mi, sagittis quis
euismod eget, sagittis in nibh. Proin elit turpis, aliquam et imperdiet
sed, volutpat eu turpis.

Pellentesque vel enim tellus, vitae egestas turpis. Praesent malesuada elit
non nisi sollicitudin non blandit lacus tincidunt. Morbi blandit urna at
lectus ornare laoreet. Suspendisse turpis diam, lobortis dictum luctus
quis, commodo at lorem. Integer lacinia convallis ultricies. Sed quis augue
neque, eu malesuada arcu. Nullam vehicula, purus vitae sagittis pulvinar,
erat eros semper massa, eu egestas nibh erat quis magna. Cras pellentesque,
nisl eu dapibus volutpat, urna augue ornare quam, quis egestas lectus nulla
a lectus.

Vivamus dictum libero in massa cursus sed vulputate eros imperdiet. Donec
lacinia, libero ac lobortis egestas, nibh dui ornare arcu, luctus porttitor
velit massa sit amet quam. Maecenas scelerisque laoreet diam, vitae congue
quam adipiscing vitae. Aliquam cursus nisl a leo convallis eleifend
fermentum massa porta. Nunc libero quam, dapibus dapibus molestie sit amet,
faucibus vel nunc.

\subsection{Praesent auctor venenatis posuere}
Sed tellus augue, molestie in pulvinar lacinia, dapibus non ipsum. Fusce
vitae mi vitae enim ullamcorper hendrerit eu malesuada est. Proin iaculis
ante sed nibh tincidunt vel interdum libero posuere. Vivamus accumsan metus
quis felis congue suscipit dapibus enim mattis. Fusce mattis tortor eget
ipsum interdum sagittis auctor id metus.

Integer diam lacus, pharetra sit amet tempor et, tristique non lorem.
Aenean auctor, nisi eu interdum fermentum, lectus massa adipiscing elit,
sed facilisis orci odio a lectus. Proin mi nibh, tempus quis porta a,
viverra quis enim. In sollicitudin egestas libero, quis viverra velit
molestie eget. Nulla rhoncus, dolor a mollis vestibulum, lacus elit semper
nisi, nec sollicitudin sem urna eu magna. Nunc sed est urna, euismod congue
mi.

\subsection{Cras vulputate ultricies venenatis}
Vivamus eros urna, sodales accumsan semper vel, lobortis sit amet mauris.
Etiam condimentum eleifend lorem, ullamcorper ornare lectus aliquet vitae.
Praesent massa enim, interdum sit amet semper et, venenatis ut elit.
Quisque faucibus, quam ac lacinia imperdiet, nulla neque elementum purus,
tempus rutrum justo massa porta sapien. Vestibulum ante ipsum primis in
faucibus orci luctus et ultrices posuere cubilia Curae; Sed ultrices
interdum mi, et rhoncus sapien rutrum sed.

Duis elit orci, molestie quis sollicitudin sed, convallis non ante.
Maecenas tincidunt condimentum justo, et ultricies leo tristique vitae.
Vestibulum quis quam non lectus dapibus eleifend a vitae nibh. Nam nibh
justo, pharetra quis iaculis consequat, elementum quis justo. Etiam mollis
lacinia lacus, nec sollicitudin urna lobortis ac. Nulla facilisi.

Proin placerat risus eleifend erat ultricies placerat. Etiam rutrum magna
nec turpis euismod consectetur. Phasellus tortor odio, lacinia imperdiet
condimentum sed, faucibus commodo erat. Phasellus sed felis id ante
placerat ultrices. Aenean tempor justo in tortor volutpat eu auctor dolor
mollis. Aenean sit amet risus urna. Morbi viverra vehicula cursus.

\subsection{Donec nibh ante, consectetur et posuere id, tempus nec arcu}
Curabitur a tellus aliquet ipsum pellentesque scelerisque. Etiam congue,
risus et volutpat rutrum, est purus dapibus leo, non cursus metus felis
eget ligula. Vivamus facilisis tristique turpis, ut pretium lectus luctus
eleifend. Fusce magna sapien, ullamcorper vitae fringilla id, euismod quis
ante.

Phasellus volutpat, nunc et pharetra semper, sem justo adipiscing mauris,
id blandit magna quam et orci. Vestibulum a erat purus, ut molestie ante.
Vestibulum ante ipsum primis in faucibus orci luctus et ultrices posuere
cubilia Curae; Proin turpis diam, consequat ut ullamcorper ut, consequat eu
orci. Sed metus risus, fringilla nec interdum vel, interdum eu nunc.
Suspendisse vel sapien orci.

\subsection{Morbi et mauris tempus purus ornare vehicula}
Mauris sit amet diam quam, eget luctus purus. Sed faucibus, risus semper
eleifend iaculis, mi turpis bibendum nisl, quis cursus nibh nisl sit amet
ipsum. Vestibulum tempor urna vitae mi auctor malesuada eget non ligula.
Nullam convallis, diam vel ultrices auctor, eros eros egestas elit, sed
accumsan arcu tortor eget leo. Vestibulum orci purus, porttitor in pharetra
eget, tincidunt eget nisl. Nullam sit amet nulla dui, facilisis vestibulum
dui.

Donec faucibus facilisis mauris ac cursus. Duis rhoncus quam sed nisi
laoreet eu scelerisque massa tincidunt. Vivamus sit amet libero nec arcu
imperdiet tempor quis non libero. Sed consequat dignissim justo. Phasellus
ullamcorper, velit quis posuere vulputate, felis erat tincidunt mauris, at
vestibulum justo lectus et turpis. Maecenas lacinia convallis euismod.
Quisque egestas fermentum sapien eu dictum. Sed nec lacus in purus dictum
consequat quis vel nisl. Fusce non urna sem. Curabitur eu diam vitae elit
accumsan blandit. Nullam fermentum nunc et leo dictum laoreet. Donec semper
varius velit vel fringilla. Vivamus eu orci nunc.

\section{Besluit van dit hoofdstuk}
Als je in dit hoofdstuk tot belangrijke resultaten of besluiten gekomen
bent, dan is het ook logisch om het hoofdstuk af te ronden met een
overzicht ervan. Voor hoofdstukken zoals de inleiding en het
literatuuroverzicht is dit niet strikt nodig.

%%% Local Variables: 
%%% mode: latex
%%% TeX-master: "masterproef"
%%% End: 

\chapter{Einde van de arbeidsovereenkomst}
\label{hoofdstuk:2}

\section{Algemene (civielrechtelijke) wijzen van beëindiging}

\section{Wederzijdse toestemming}

\section{Overmacht}

\section{Ontbindend beding}

\section{Afloop van de termijn of voltooiing van het afgesproken werk}

\section{Overlijden van een van de partijen}

\section{Arbeidsrechtelijke wijzen van beëindiging: 'ontslag' en 'opzeg'}

\subsection{Algemeen}

\subsection{De opzeg}

\subsubsection{Algemeen}

\subsubsection{Wijze van kennisgeving van de opzeg}

\paragraph{Opzeg uitgaande van de werknemer}

\paragraph{Opzeg uitgaande van de werkgever}

\subsubsection{Inhoud}

\subsubsection{Nietigheid van de opzeg}

\subsubsection{Opzegtermijn}

\paragraph{Begrip - Wat voorafging aan de Wet Eenheidsstatuut}

\paragraph{Aanvang van de termijn}

\paragraph{Belang van de anciënniteit voor de duur van de opzegtermijn}

\paragraph{Duur van de opzegtermijn: algemeen}

\paragraph{Opzegtermijn: algemene regels}

\paragraph{Uitzonderingen op de opzegtermijnen}

\paragraph{Geen afwijkingen in een cao gesloten in een paritair (sub)comité}

\paragraph{Opzegtermijn in geval van pensioen}

\paragraph{Schorsing van de opzegtermijn}

\subparagraphsection{Algemeen}

\subparagraphsection{Schorsingsoorzaken en invloed op de opzegtermijn}

\paragraph{Recht op afwezigheid tijdens de opzegtermijn: sollicitatieverlof}

\paragraph{De opzegvergoeding}

\subparagraphsection{Algemeen}

\subparagraphsection{Loon als berekeningsbasis voor de opzegvergoeding}

\section{Ontslag wegens dringende reden}

\subsection{Algemeen}

\subsubsection{Begrip dringende reden}

\subsubsection{Voorwaarden}

\paragraph{"Een ernstige tekortkoming"}

\paragraph{"Die elke professionele samenwerking onmogelijk maakt"}

\paragraph{"De samenwerking moet onmiddelijk en definitief onmogelijk gemaakt zijn"}

\subsection{Termijnen}

\subsubsection{Termijn om tot ontslag over te gaan}

\subsubsection{Termijn en kennisgeving van de motieven}

\subsubsection{Bewijs}

\subsubsection{Vergoedingsregeling}

\subsubsection{Voorbeelden uit de rechtspraak}

\section{Motivering van het ontslag en kennelijk onredelijk ontslag}

\subsection{Algemeen}

\subsection{Motiveringsplicht}

\subsection{Kennelijk onredelijk ontslag}

\subsubsection{Inhoud}

\subsubsection{Uitzonderingen}

\subsection{Willekeurig ontslag van arbeiders}

\section{Onregelmatige beëindiging of verbreking van de arbeidsovereenkomst}

\subsection{Algemeen}

\subsection{Opzegvergoeding}

\section{Bescherming tegen ontslag}

\section{Rechtsmisbruik bij ontslag}

\subsection{Algemeen}

\subsection{Rechtsmisbruik bij ontslag}

\section{Het ontslag van een arbeidsongeschikte werknemer}




%%% Local Variables: 
%%% mode: latex
%%% TeX-master: "masterproef"
%%% End: 

\chapter{De reglementering van de arbeid}
\label{hoofdstuk:2}

\section{Het arbeidsreglement}

\subsection{Wat is een arbeidsreglement?}

\subsection{Vorm en inhoud}

\subsubsection{Verplichte bepalingen volgens de wet op het arbeidsreglement}

\paragraph{De arbeidstijdregeling}

\paragraph{De arbeid en het loon}

\paragraph{Het einde van de overeenkomst}

\paragraph{De rechten en plichten van het toezichthoudend personeel}

\paragraph{Straffen, boeten en verhaalprocedure}

\paragraph{De veiligheid en de eerste hulp bij ongevallen of EHBO}

\paragraph{Het medisch getuigschrift (attest) bij ziekte}

\paragraph{De medische tussenkomst bij arbeidsongeval}

\paragraph{De organen van medezeggenschap en de syndicale afvaardiging}

\paragraph{De jaarlijkse vakantie en de betaalde feestdagen}

\paragraph{De inspectie}

\paragraph{Het onthaal en de aanpassing van de werknemers}

\paragraph{Collectieve arbeidsovereenkomsten}

\paragraph{Elektronische arbeidsovereenkomst en andere sociale documenten}

\subsubsection{Verplichte bepalingen volgens andere wetten en besluiten}

\subsubsection{Vrije bepalingen}

%%% Local Variables: 
%%% mode: latex
%%% TeX-master: "masterproef"
%%% End: 




% ... en zo verder tot
\chapter{Het laatste hoofdstuk}
\label{hoofdstuk:n}
Een hoofdstuk behandelt een samenhangend geheel dat min of meer op zichzelf
staat. Het is dan ook logisch dat het begint met een inleiding, namelijk
het gedeelte van de tekst dat je nu aan het lezen bent.

\section{Eerste onderwerp in dit hoofdstuk}
De inleidende informatie van dit onderwerp.

\subsection{Een item}
De bijbehorende tekst. Denk eraan om de paragrafen lang genoeg te maken en
de zinnen niet te lang.

Een paragraaf omvat een gedachtengang en bevat dus steeds een paar zinnen.
Een paragraaf die maar \'e\'en lijn lang is, is dus uit den boze.

\section{Tweede onderwerp in dit hoofdstuk}
Er zijn in een hoofdstuk verschillende onderwerpen. We zullen nu
veronderstellen dat dit het laatste onderwerp is.

\section{Besluit van dit hoofdstuk}
Als je in dit hoofdstuk tot belangrijke resultaten of besluiten gekomen
bent, dan is het ook logisch om het hoofdstuk af te ronden met een
overzicht ervan. Voor hoofdstukken zoals de inleiding en het
literatuuroverzicht is dit niet strikt nodig.



%%% Local Variables: 
%%% mode: latex
%%% TeX-master: "masterproef"
%%% End: 


\part{ICT-RECHT}

%\chapter{Algemene bepalingen van de arbeidsovereenkomstenwet van          3 juli 1978}
\label{hoofdstuk:1}

%\setcounter{section}{-1}

\section{Toepassingsgebied van de Arbeidsovereenkomst}

\subsection{Wat is een arbeidsovereenkomst?}

\subsection{Het ondergeschikte verband}

\subsubsection{Algemeen}

\subsubsection{Schijnzelfstandigen - Pseudowerknemers}

\subsubsection{Samengaan van een arbeidsovereenkomst met een aannemings- of zelfstandigenovereenkomst}

\subsection{Toepassingsgebied van de Arbeidsovereenkomstenwet}

\section{Het aangaan van de arbeidsovereenkomst}

\subsection{Geldigheidsvereisten}

\subsubsection{Bekwaamheid}

\subsubsection{Toestemming}

\paragraph{Geweld}

\paragraph{Dwaling}

\paragraph{Bedrog}

\subsubsection{Voorwerp en oorzaak}

\subsubsection{Gevolgen}

\subsection{Vormvereisten van de arbeidsovereenkomst}

\section{Soorten arbeidsovereenkomsten}

\subsection{Soorten naargelang de beoogde arbeid}

\subsubsection{Werklieden}

\subsubsection{Bedienden}

\subsubsection{Handelsvertegenwoordigers}

\subsubsection{Dienstboden}

\subsubsection{Studenten}

\subsubsection{Huisarbeid}

\subsection{Soorten naargelang de duur van de overeenkomst}

\subsubsection{De arbeidsovereenkomst voor onbepaalde tijd}

\subsubsection{De arbeidsovereenkomst voor een bepaalde tijd}

\paragraph{Onderscheid tussen overeenkomsten van onbepaalde en bepaalde tijd}

\paragraph{Opeenvolgende arbeidsovereenkomsten voor een bepaalde tijd}

\paragraph{Afwijking van het principiële verbod om opeenvolgende arbeidsovereenkomsten voor een bepaalde tijd te sluiten}

\paragraph{Algemeen beginsel van non-discriminatie ten gunste van werknemers met een overeenkomst voor bepaalde tijd}

\subsubsection{De arbeidsovereenkomst voor duidelijk omschreven werk}

\subsubsection{De vervangingsovereenkomst}

\subsection{Soorten naargelang de omvang van het werk}

\subsubsection{Voltijdse arbeidsovereenkomst}

\subsubsection{Arbeidsovereenkomst voor deeltijdse arbeid}

\paragraph{Geschrift}

\paragraph{De arbeidsregeling}

\paragraph{Het werkrooster}

\paragraph{Minimumgrenen van de arbeidsduur}

\paragraph{Bekendmaking van de variabele werkroosters}

\paragraph{Voorrang bij het overstappen naar een voltijdse betrekking}

\paragraph{Non-discriminatie ten gunste van deeltijdse werknemers}

\paragraph{Rechten deeltijdse werknemer versus voltijdse werknemer}


\subsection{Soorten naargelang de wijze waarop de overeenkomst werd gesloten}

\subsubsection{Mondelinge en schriftelijke arbeidsovereenkomsten}

\subsubsection{Elektronische arbeidsovereenkomsten}

%%% Local Variables: 
%%% mode: latex
%%% TeX-master: "masterproef"
%%% End: 


%\chapter{Besluit}
\label{besluit}
De masterproeftekst wordt afgesloten met een hoofdstuk waarin alle
besluiten nog eens samengevat worden. Dit is ook de plaats voor suggesties
naar het verder gebruik van de resultaten, zowel industri"ele toepassingen
als verder onderzoek.

\lipsum[1-7]

%%% Local Variables: 
%%% mode: latex
%%% TeX-master: "masterproef"
%%% End: 


%---------
\aliaspagestyle{chapter}{nohead}

% Indien er bijlagen zijn:
\appendixpage*          % indien gewenst
\appendix

%allemaal new stuff

%\appendixpage          % indien gewenst
%\appendixtableofcontents
%end allemaal new stuff

%nieuw
%\import{bijlagen/}{bijlage1_opdrachtenlijst.tex}

\chapter{De eerste bijlage}
\label{app:A}
In de bijlagen vindt men de data terug die nuttig kunnen zijn voor de
lezer, maar die niet essentieel zijn om het betoog in de normale tekst te
kunnen volgen. Voorbeelden hiervan zijn bronbestanden,
configuratie-informatie, langdradige wiskundige afleidingen, enz.

In een bijlage kunnen natuurlijk ook verdere onderverdelingen voorkomen,
evenals figuren en referenties\cite{h2g2}.

\section{Meer lorem}
\lipsum[50]

\subsection{Lorem 15--17}
\lipsum[15-17]

\subsection{Lorem 18--19}
\lipsum[18-19]

\section{Lorem 51}
\lipsum[51]

%%% Local Variables: 
%%% mode: latex
%%% TeX-master: "masterproef"
%%% End: 

% ... en zo verder tot
\chapter{De laatste bijlage}
\label{app:n}
In de bijlagen vindt men de data terug die nuttig kunnen zijn voor de
lezer, maar die niet essentieel zijn om het betoog in de normale tekst te
kunnen volgen. Voorbeelden hiervan zijn bronbestanden,
configuratie-informatie, langdradige wiskundige afleidingen, enz.

\section{Lorem 20-24}
\lipsum[20-24]

\section{Lorem 25-27}
\lipsum[25-27]

%%% Local Variables: 
%%% mode: latex
%%% TeX-master: "masterproef"
%%% End: 


%\includepdf[pages=-]{bijlagen/website_audit_filip_eline_belgian_interns.pdf}


\backmatter
% Na de bijlagen plaatst men nog de bibliografie.

%\pagestyle{afterbackmatterf}

%\chapter{Bibliografie}


%for bibLatex
%\printbibliography[heading=bibintoc,title={Whole bibliography}] %Prints the entire bibliography with the titel "Whole bibliography"

%hiermee worden alle entries getoond, zelfs als ze niet vermeld worden in de tekst.
%\nocite{*}

%Filters bibliography
%\printbibliography
\printbibliography[title={Referenties}]
\pagestyle{afterbackmatterf}

%\printbibliography[heading=subbibintoc,type=online,title={Sites only}]
%\todo[inline]{Gebruikte bronnen en sites nog als bibliografie toevoegen achteraan.}

%\chapter{Glossaries}

%\glsaddall
%\printglossaries

\setglossarystyle{altlisthypergroup}

\printglossary[title=Termen,toctitle=Lijst van termen]

%\clearpage

\setglossarystyle{mcolindex}

\printglossary[type=\acronymtype]
%\printglossary[type=\acronymtype,title=Abbreviations]

\chapter{Index}

\printindex

\end{document}

%%% Local Variables: 
%%% mode: latex
%%% TeX-master: t
%%% End: 
